\subsection{関連概念}

\subsubsection{会計}

会計は、組織の資金がどのように使われ、どのような成果を生んだかを測定し、利害関係者へ伝える活動である。営利組織では投資収益率、非営利組織や政府では持続可能性や資金使用の妥当性に関係する。会計システムは、見積りに使う過去データの重要な源泉でもある。

\subsubsection{費用と原価計算}

費用とは、何かを得るために使われ、別用途には使えなくなった資源である。経済学的には、ある選択によって諦めた代替案も費用である。

埋没費用は、すでに発生して回収できない費用である。機会費用は、ある案を選ぶことで諦めた別案の価値である。原価計算は、開発、製造、配布、再作業、保守などにかかる費用を把握する活動である。

総所有費用は、取得、導入、運用、保守、廃止まで含めた総費用である。ソフトウェアでは初期開発費だけを見ると誤る。運用、保守、教育、移行、障害対応、セキュリティ対応、廃止費用まで考える必要がある。

\subsubsection{財務}

財務は、資源を調達、配分、管理、投資する活動である。時間、資金、リスクの関係を扱い、組織の戦略、資金調達、ポートフォリオ、キャッシュフロー、財務業績の測定に関わる。ソフトウェアは現代の組織の財務に大きく影響するため、技術者も財務の基本を理解する必要がある。

\subsubsection{統制}

統制とは、実績を測定し、計画との差を把握し、必要な修正を行うことである。費用統制では、計画費用と実費用の差を検出する。ソフトウェア工学では、プロセスや製品の状態を測定し、組織目標と整合するよう改善することに関係する。

\subsubsection{効率と有効性}

効率は「正しく行う」ことである。少ない資源で期待される成果を出すことに関係する。有効性は「正しいことを行う」ことである。定義された目的を達成したかを見る。

\begin{notebox}{用語の見分け方}
効率だけを高めても、そもそも価値のない機能を速く作っているなら意味は小さい。有効性だけを見ても、過剰な資源を使っていれば持続しにくい。生産性は、効率と有効性を価値の観点で結び付ける。
\end{notebox}

\subsubsection{生産性}

生産性は、入力に対する出力の比率である。出力は提供された価値であり、入力は費やされた人員、時間、資金などである。生産性を高めるには、より高い価値をより少ない資源で生む必要がある。

ソフトウェア開発では、再作業が最大の資源消費になることが多い。欠陥を早期に見つける、欠陥の影響拡大を抑える、欠陥を最初から防ぐといった品質改善は、生産性向上に直結する。

\subsubsection{製品またはサービス}

製品は、入力を変換して作られる有形または経済的な成果物である。サービスは、コンサルティングのような無形の資源提供である。ソフトウェアは、社内 IT ソリューション、外販アプリケーション、組込み部品、クラウドサービスなど、製品とサービスの両方の性質を持つことがある。

\subsubsection{プロジェクト}

プロジェクトは、独自の製品、サービス、結果を作るための一時的な取り組みである。ソフトウェア製品のライフサイクル中には、企画、初期開発、機能追加、保守、移行、廃止など、複数のプロジェクトが含まれる。

\subsubsection{プログラム}

プログラムは、関連する複数のプロジェクトや活動を調整して管理する単位である。個々のプロジェクトを独立に管理するだけでは得られない便益を得るために使う。

\subsubsection{ポートフォリオ}

ポートフォリオは、戦略目標を達成するために、プロジェクト、プログラム、サブポートフォリオ、運用活動をまとめて管理する単位である。あるプロジェクトへ資源を配分することは、他のプロジェクトへ同じ資源を使えないことを意味する。したがって、ポートフォリオ視点は機会費用の理解に役立つ。

\subsubsection{製品ライフサイクル}

ソフトウェア製品ライフサイクルは、製品やサービスを定義、構築、運用、保守、廃止するすべての活動を含む。初期開発だけではなく、運用、保守、廃止の方が長期的には多くの資源を消費することが多い。したがって、経済性は初期開発費だけでは判断できない。

\subsubsection{プロジェクトライフサイクル}

プロジェクトライフサイクルは、開始、計画、実行、監視と統制、終結などの活動を含む。アジャイル開発では、複数の反復が成果物を増分的に作る。製品ライフサイクルには、複数の開発ライフサイクルが含まれることがある。

\subsubsection{価格と価格設定}

価格は、財やサービスと引き換えに支払われるものである。価格設定は、企業が製品やサービスの対価として何を受け取るかを決める活動である。製造費、競争、市場状態、品質、販売キャンペーン、数量割引、サービス内容などが影響する。

\subsubsection{優先順位付け}

優先順位付けは、共通基準に基づいて代替案を並べ、最も価値の高い順に扱うことである。要求の優先順位付けは、限られた納期、予算、人員、技術制約の中で、顧客へ最も価値を届けるために使われる。

\par\smallskip
\begin{center}
\centering
\IfFileExists{assets/generated_figures/ch15/fig-ch15-09.png}{\includegraphics[width=0.96\linewidth]{assets/generated_figures/ch15/fig-ch15-09.png}}{\fbox{\parbox[c][48mm][c]{0.9\linewidth}{\centering 画像生成待ち\\SDLC と SPLC の違いを示す図}}}
\par\smallskip
\refstepcounter{figure}\label{fig:ch15-09}図 \thefigure: SDLC と SPLC の違いを示す図
\end{center}
図 \thefigure は、SDLC と SPLC の違いを示す図を視覚的に整理したものである。
\par\smallskip
